% !TEX program = xelatex
% =============================================================================
%  全国大学生数学建模竞赛（CUMCM）LaTeX 模板
%  本项目的 2026 CUMCM LaTeX 独立模板；本次规范改造不影响 Typst 模板。
%
%  字体依赖（正文宋体 / 标题黑体 / 强调楷体）：
%    - 规范字体：SimSun（宋体）、SimHei（黑体）、KaiTi（楷体）。
%    - 本模板默认使用 ctex 的 fontset=mac，在 macOS 上自动选用同族等价字体
%      Songti SC / Heiti SC / Kaiti SC，确保开箱即用。
%    - 若系统已安装 SimSun/SimHei/KaiTi（如 Windows 或手动安装），可将下方
%      \documentclass 选项 fontset=mac 改为 fontset=windows，即可使用规范字体。
%    - 西文正文：Times New Roman；代码等宽：Menlo（缺失时回退到默认字体）。
%  编译方式：xelatex main.tex（建议运行两遍以更新交叉引用）
% =============================================================================
\documentclass[fontset=mac, 12pt, a4paper]{ctexart}

% ---------- 页面：A4，上下左右 2.5cm ----------
\usepackage[a4paper, top=2.5cm, bottom=2.5cm, left=2.5cm, right=2.5cm]{geometry}

% ---------- 数学 ----------
\usepackage{amsmath}

% ---------- 字体（fontspec + ctex，xelatex 编译） ----------
% 中文宋体/黑体/楷体由 ctex 的 fontset=mac 自动设置（Songti SC / Heiti SC / Kaiti SC），
% 并自动定义 \songti / \heiti / \kaishu 等字体命令。
% 西文正文与代码等宽字体：
\IfFontExistsTF{Times New Roman}{\setmainfont{Times New Roman}}{}
\IfFontExistsTF{Menlo}{\setmonofont{Menlo}}{}

% ---------- 正文：12.05pt，首行缩进 2em，两端对齐，行距 0.72em ----------
% 12pt 字号近似 12.05pt；0.72em 行距 => 行高 = 1.72 * 字号，\linespread 因子 ≈ 1.43
\linespread{1.43}
\setlength{\parindent}{2em}

% ---------- 标题编号（\ctexset 配置） ----------
% 一级标题用中文数字“一、二、三、”，二级及以下用 1.1 / 1.1.1
\ctexset{
  section/number       = \chinese{section},
  subsection/number    = \arabic{section}.\arabic{subsection},
  subsubsection/number = \arabic{section}.\arabic{subsection}.\arabic{subsubsection},
}

% ---------- 标题格式（titlesec 配置） ----------
\usepackage{titlesec}
\titleformat{\section}
  {\centering\fontsize{17.3pt}{20.8pt}\heiti\bfseries}
  {\chinese{section}、}{1em}{}
\titleformat{\subsection}
  {\fontsize{14.45pt}{17.4pt}\heiti\bfseries}
  {\arabic{section}.\arabic{subsection}}{1em}{}
\titleformat{\subsubsection}
  {\fontsize{12.05pt}{14.5pt}\heiti\bfseries}
  {\arabic{section}.\arabic{subsection}.\arabic{subsubsection}}{1em}{}
\titlespacing{\section}       {0pt}{1.25em}{0.82em}
\titlespacing{\subsection}    {0pt}{1.15em}{0.55em}
\titlespacing{\subsubsection} {0pt}{1.15em}{0.55em}

% ---------- 表格与代码 ----------
\usepackage{booktabs}
\usepackage{array}
\usepackage{xcolor}
\usepackage{listings}
\lstset{
  basicstyle=\ttfamily\small,
  backgroundcolor=\color{black!3},
  frame=single,
  framesep=6pt,
  rulecolor=\color{black!30},
  framerule=0.8pt,
  breaklines=true,
  showstringspaces=false,
  columns=fullflexible,
  keepspaces=true,
}

% ---------- 页码：底部居中 ----------
\pagestyle{plain}

% ---------- 自定义命令（对应 Typst 函数） ----------
% \papertitle{标题}：居中 17.3pt 加粗标题
\newcommand{\papertitle}[1]{%
  \thispagestyle{plain}%
  \begin{center}
    {\fontsize{17.3pt}{22pt}\bfseries #1}
  \end{center}
  \vspace{1em}
}
% \abstractcn{摘要正文}{关键词}：居中“摘要”标题（14pt 加粗）+ 正文 + “关键字：”加粗 + 关键词 + 换页
\newcommand{\abstractcn}[2]{%
  \thispagestyle{plain}%
  \begin{center}{\fontsize{14pt}{17pt}\bfseries 摘要}\end{center}
  #1\par
  \vspace{1em}
  {\heiti\bfseries 关键字：}#2\par
  \newpage
}
% \referencescn：参考文献标题由 thebibliography 自动生成，本宏不重复设置标题
\newcommand{\referencescn}{%
  %%%%%%%%%%%%%%%%%%%%%%%%%%%%%%%%%%%%%%%%%%%%%%%%%%%%%%%%%%%%%
% ==================== 参考文献 ====================
% ✓ 自查清单：
%   □ 是否只收录真实存在且在正文中使用 \cite{} 引用的资料？
%   □ 是否按正文首次引用顺序编号？
%   □ 作者、题名、年份、卷期、页码或网址是否完整且准确？
%   □ 是否删除未在正文引用的占位条目？
%   □ 是否未将 ChatGPT、DeepSeek、Claude、Copilot 等 AI 工具列为参考文献？
%
% 注意：thebibliography 会自动生成“参考文献”标题，本文件不要再写 \section*。

\begin{thebibliography}{99}
  \bibitem{ref-journal} 【作者】. 【文章题名】[J]. 【期刊名称】, 【年份】, 【卷】(【期】): 【起止页码】.
  \bibitem{ref-book} 【作者】. 【书名】[M]. 【出版地】: 【出版社】, 【出版年】.
  \bibitem{ref-thesis} 【作者】. 【学位论文题名】[D]. 【城市】: 【授予单位】, 【年份】.
  \bibitem{ref-conference} 【作者】. 【会议论文题名】[C]//【论文集或会议名称】. 【出版地】: 【出版者】, 【年份】: 【起止页码】.
  \bibitem{ref-online} 【作者或机构】. 【网页或数据资料题名】[EB/OL]. 【发布日期】[【引用日期】]. 【网址】.
\end{thebibliography}
%
}
% \appendixcn：附录使用 A、B、C、D 编号
\newcommand{\appendixcn}{%
  \appendix
  \titleformat{\section}
    {\centering\fontsize{17.3pt}{20.8pt}\heiti\bfseries}
    {附录~\Alph{section}\quad}{0pt}{}
  \titleformat{\subsection}
    {\fontsize{14.45pt}{17.4pt}\heiti\bfseries}
    {\Alph{section}.\arabic{subsection}}{1em}{}
  \titleformat{\subsubsection}
    {\fontsize{12.05pt}{14.5pt}\heiti\bfseries}
    {\Alph{section}.\arabic{subsection}.\arabic{subsubsection}}{1em}{}
  %%%%%%%%%%%%%%%%%%%%%%%%%%%%%%%%%%%%%%%%%%%%%%%%%%%%%%%%%%%%%
% ==================== 附录 ====================
% 附录依据真实 delivery_manifest 和支撑材料生成；代码部分只选择主要建模、求解和有方法价值的验证程序。
% 最终提交前必须删除全部【占位内容】，并确认文件名、代码与支撑材料压缩包完全一致。

\section{支撑材料文件列表}

支撑材料中包含用于支撑模型、结果和结论的必要文件，具体如表~\ref{tab:supporting-files} 所示。
文件名和功能描述必须根据真实交付文件填写，不得预置不存在的数据、程序或结果。
若确实没有支撑材料，应删除下表并明确写明“本论文没有支撑材料”。

\begin{table}[htbp]
  \centering
  \caption{支撑材料文件列表}
  \begin{tabular}{p{0.32\textwidth}p{0.58\textwidth}}
    \toprule
    \textbf{文件名} & \textbf{功能描述} \\
    \midrule
    【真实文件名】 & 【说明该文件对应的问题、输入、计算或结果】 \\
    【真实文件名】 & 【说明该文件在论文中的支撑作用】 \\
    \ifcumcmaiused
      AI工具使用详情.pdf & 记录AI工具名称、使用环节、主要提示方式、采纳修改与人工核验情况 \\
    \fi
    \bottomrule
  \end{tabular}
  \label{tab:supporting-files}
\end{table}

\clearpage
\section{计算环境与复现说明}

本文所用计算环境和求解工具如表~\ref{tab:computing-environment} 所示。
版本、硬件和运行命令必须来自真实环境；若某类工具未使用，应删除对应行。

\begin{table}[htbp]
  \centering
  \caption{计算环境与求解工具}
  \begin{tabular}{p{0.25\textwidth}p{0.65\textwidth}}
    \toprule
    \textbf{项目} & \textbf{真实配置} \\
    \midrule
    操作系统与硬件 & 【操作系统、处理器、内存等】 \\
    编程语言或软件 & 【名称及准确版本】 \\
    主要依赖库 & 【库名称及准确版本】 \\
    优化器或求解器 & 【名称、版本及关键设置；未使用则删除】 \\
    随机性控制 & 【随机种子、重复运行次数或不适用】 \\
    运行方式 & 【入口文件、命令和必要参数】 \\
    \bottomrule
  \end{tabular}
  \label{tab:computing-environment}
\end{table}

为复现论文结果，应依次执行【真实运行顺序】，输入文件为【实际文件】，主要输出为【实际文件】。

\clearpage
\section{主要建模与求解源程序}

本附录根据支撑材料中的真实程序和 \texttt{delivery\_manifest.json}，选择能够体现数学模型、核心算法、求解过程以及重要验证方法的完整必要代码。不得把整个工程目录机械复制到论文，也不能只给出伪代码、空框架或文件名。

仅负责参数解析和函数转发的入口程序、简单字段检查以及纯绘图脚本，默认不收入本附录。这里的纯绘图脚本是指只负责字体、坐标轴、图例、标注、子图排版和保存图片的程序；其完整版本仍保留在电子支撑材料的 \texttt{code/plots/} 中。

以下问题小节和程序数量必须按实际赛题与主要算法动态增删。若确实没有使用程序，应删除本节占位内容并明确写明“本论文没有用到程序”。

\subsection{支撑材料中的主要源程序列表}

\begin{table}[htbp]
  \centering
  \caption{论文附录选取的主要源程序}
  \begin{tabular}{p{0.16\textwidth}p{0.34\textwidth}p{0.40\textwidth}}
    \toprule
    \textbf{对应问题} & \textbf{真实程序文件} & \textbf{选取理由与程序作用} \\
    \midrule
    【问题编号】 & 【支撑材料中的相对路径】 & 【说明其体现的模型、算法或重要验证方法】 \\
    \bottomrule
  \end{tabular}
  \label{tab:appendix-main-code}
\end{table}

\subsection{问题一主要源程序}

\noindent\textbf{选取文件：}【问题一实际主要模型、求解或重要验证程序】

\noindent\textbf{选取理由：}【说明该代码在问题一建模与求解中的核心作用】

\noindent 【使用 \verb|\lstinputlisting| 插入被选中的完整必要代码；不存在的模块不得保留。】
% \lstinputlisting[language=Python]{../../code/solvers/problem1/model.py}
% \lstinputlisting[language=Python]{../../code/solvers/problem1/solver.py}

\subsection{问题二主要源程序}

\noindent 【按照问题一格式，根据问题二实际主要程序填写。】

\subsection{问题三主要源程序}

\noindent 【按照问题一格式，根据问题三实际主要程序填写。】

\subsection{问题四主要源程序}

\noindent 【按照问题一格式，根据问题四实际主要程序填写；若题面不是四问则动态调整。】

\clearpage
\section{补充推导、数据与计算结果}

本节用于放置正文未完整展示但支撑结论所必需的长篇公式推导、长表和补充图表。
所有内容必须能够追溯至真实数据或求解记录；若没有补充内容，应删除本节。

\begin{table}[htbp]
  \centering
  \caption{【补充数据或计算结果名称】}
  \begin{tabular}{cccc}
    \toprule
    \textbf{【字段一】} & \textbf{【字段二】} & \textbf{【字段三】} & \textbf{【字段四】} \\
    \midrule
    【真实数据】 & 【真实数据】 & 【真实数据】 & 【真实数据】 \\
    \bottomrule
  \end{tabular}
  \label{tab:supplementary-results}
\end{table}
%
}
% \threelinetable{标题}{列定义}{表头}{表格内容}：三线表，标题在表格上方居中 10.5pt 黑体加粗
% 调用示例：\threelinetable{表 1 示例}{ccc}{A & B & C}{1 & 2 & 3 \\ 4 & 5 & 6}
\newcommand{\threelinetable}[4]{%
  \begin{center}
    {\fontsize{10.5pt}{12.6pt}\heiti\bfseries #1}\par
    \vspace{0.6em}
    \begin{tabular}{#2}
      \toprule
      #3 \\
      \midrule
      #4 \\
      \bottomrule
    \end{tabular}
  \end{center}
}

\begin{document}

% 电子版第一页为摘要页，页码从 1 开始
\setcounter{page}{1}
\papertitle{论文标题}

\abstractcn{%
  中文摘要内容：问题概述 + 每个子问题的方法和数值结果 + 结论%
}{%
  关键词1 \quad 关键词2 \quad 关键词3%
}

\input{sections/1_restatement}
%%%%%%%%%%%%%%%%%%%%%%%%%%%%%%%%%%%%%%%%%%%%%%%%%%%%%%%%%%%%%
% ==================== 第二章：问题分析 ====================
% BZD 数模社制作——国赛 B 题四问通用模板
% 使用时请根据题目和附件逐一替换【】中的内容，并删除不适用的表述。
%
% ✓ 自查清单：
%   □ 是否说明了每个问题的任务、研究意义和核心困难？
%   □ 是否判断了各问题的数学类型，并说明方法选择依据？
%   □ 是否分析了附件数据的内容、结构、质量和适用范围？
%   □ 是否说明了数据分析、清洗和预处理方法及其理由？
%   □ 是否明确了题目要求的输出、约束和评价标准？
%   □ 是否仅在确有必要时设置主模型与对比模型？
%   □ 是否说明了四个问题之间的数据或逻辑联系？
%   □ 是否规划了误差分析、灵敏度分析或稳健性检验？
%%%%%%%%%%%%%%%%%%%%%%%%%%%%%%%%%%%%%%%%%%%%%%%%%%%%%%%%%%%%%

\section{问题分析}

\subsection{问题一的分析}

问题一要求在【研究背景或已知条件】下完成【问题一的具体任务】，其研究意义在于【说明该结果对认识研究对象、揭示基本规律或解决实际问题的作用】。本问需要重点回答【核心问题】，并为后续问题提供【参数、特征、状态量、中间结果或基础结论】。

从数学建模角度看，问题一属于【机理分析、统计分析、预测、分类、优化或其他数学问题类型】，核心在于【需要刻画的关系、规律或决策目标】。解决此类问题通常可以采用【候选方法类别】，但不同方法对【样本规模、变量关系、分布特征、约束条件或先验假设】的要求不同，因此需要结合题目条件和附件数据进行选择，而不能仅依据模型复杂程度确定方案。

附件【附件名称或编号】主要包含【数据对象、变量和时间或空间范围】。应首先通过【描述性统计、分布分析、相关性分析、时序或空间特征分析】认识数据，并检查【缺失值、异常值、重复记录、量纲差异、统计口径或变量相关性】。针对发现的问题，采用【缺失值处理、异常值识别、数据变换、标准化或特征构造方法】，其理由是【说明该处理对保持数据真实性、满足模型条件或提高结果可靠性的作用】。

题目要求输出【问题一的结果形式】，并满足【精度、范围、物理意义或其他约束】，可利用【误差指标、拟合优度、约束回代或其他评价标准】检验结果。基于上述分析，首先建立【主模型名称或类型】刻画【主要关系】，必要时再建立【对比或备选模型】进行验证，并依据【比较指标】选择更合适的结果。问题一得到的【关键输出】将作为问题二中【对应变量或条件】的输入。

\subsection{问题二的分析}

问题二在问题一的【关键结果或基础模型】上进一步加入【新增条件、影响因素或研究范围】，要求完成【问题二的具体任务】。研究该问题有助于【说明其对模型深化、实际应用或决策改进的意义】，其核心是分析【新增因素】如何影响【研究对象或目标结果】。

从数学建模角度看，本问属于【参数估计、关系分析、预测、优化、评价或其他数学问题类型】，与问题一相比，新增难点在于【非线性、动态性、耦合关系、多重约束或不确定性】。因此，需要在问题一方法的基础上判断【原模型是否仍然适用】，并根据新增条件选择【模型扩展、参数修正、分层建模、组合建模或其他处理思路】。

本问使用的数据由【问题一的中间结果】与附件【附件名称或编号】中的【新增数据】共同构成。建模前需要检查二者在【时间范围、空间尺度、量纲、统计口径和对象标识】上的一致性，并通过【数据匹配、口径统一、相关性或因果性分析、特征筛选】识别真正影响本问结果的因素。若数据存在【样本不平衡、信息缺失、共线性或噪声干扰】，则采用【相应处理方法】，以避免【说明可能产生的偏差或失真】。

题目要求给出【问题二的结果形式】，并重点考察【准确性、稳定性、可解释性、计算效率或可实施性】。据此，可在问题一模型基础上构建【改进后的主模型】，并在有比较价值时选用【备选模型或基准方法】，利用【评价指标、交叉验证、对比实验或约束检验】比较结果。最终得到的【参数、预测结果、判定结果或方案】将为问题三中的【综合分析或最终决策任务】提供依据。

\subsection{问题三的分析}

问题三综合问题一和问题二得到的【基础结果与改进结果】，在【新的情景、约束或应用要求】下完成【问题三的具体任务】。本问是前述研究的进一步拓展，其意义在于【说明该结果对刻画复杂规律、完善模型或支持后续综合决策的价值】。

从数学建模角度看，本问属于【综合决策、多目标优化、情景分析、机理推演或其他数学问题类型】，核心难点是同时处理【多个目标、约束耦合、不确定因素或结果冲突】。因此，需要明确【决策变量或待求量】、【目标或评价准则】及【主要约束】，并分析问题一、问题二的结果如何转化为本问的【模型参数、边界条件或评价依据】。

本问需要整合附件【附件名称或编号】及前两问生成的【中间数据或模型输出】。应检查数据传递过程中的【误差累积、单位一致性、时间或空间对齐以及适用范围】，并通过【数据融合、情景划分、参数标定或不确定性描述方法】形成可用于求解的统一数据基础。对于附件无法直接给出的【参数或条件】，只能依据【明确的数据来源、题设条件或可验证的计算关系】确定，不得为了得到理想结果而人为设定。

题目要求输出【问题三的预测值、判定结果、优化方案或中间结论】，并满足【可行性、稳定性、精度或实际约束】。基于上述分析，可建立【综合主模型】完成求解；若确有必要，再以【基准模型、替代算法或不同情景】进行对比，并通过【误差分析、约束回代或对比实验】检验结果。问题三得到的【关键参数、情景结果或候选方案】将作为问题四中【综合决策或推广任务】的基础。

\subsection{问题四的分析}

问题四在前三问的【模型、参数、规律或候选方案】基础上，进一步考虑【更复杂的实际情景、新增限制或推广要求】，要求完成【问题四的具体任务】。本问承担全文的综合应用任务，其研究意义在于【说明最终结果对实际决策、系统改进、方案推广或管理建议的价值】。

从数学建模角度看，本问属于【综合决策、多目标优化、系统评价、情景推演或其他数学问题类型】，核心难点是协调【多个研究对象、多个目标、耦合约束、动态变化或不确定因素】。因此，需要明确【最终决策变量或待求量】、【总体目标或评价准则】和【必须满足的约束】，并说明前三问的结果如何转化为本问的【输入参数、边界条件、候选集合或评价依据】。

本问需要整合全部附件数据以及前三问生成的【中间数据、模型参数和求解结果】。应重点检查不同来源数据之间的【对象对应关系、时间或空间尺度、单位与统计口径】，并评估前序结果传递产生的【误差累积和不确定性】。针对【不同情景、参数波动或数据扰动】，可采用【情景划分、数据融合、不确定性描述或参数扰动方法】构造可比较的输入条件，从而保证最终结论具有明确的数据依据。

题目最终要求输出【问题四的综合结果、最优方案、实施策略或建议】，并满足【可行性、稳定性、经济性、时效性或其他实际约束】。基于上述分析，可建立【最终主模型】统一求解；若存在多种可行方法，则选用【基准模型、替代算法或不同情景】进行比较，并依据【综合评价指标】确定最终结果。最后通过【灵敏度分析、稳健性检验、误差分析、约束回代或对比实验】说明结论在【关键参数变化、数据扰动或应用场景改变】下是否仍然成立，并据此给出【与题目要求一致的最终结论或建议】。

%%%%%%%%%%%%%%%%%%%%%%%%%%%%%%%%%%%%%%%%%%%%%%%%%%%%%%%%%%%%%
% ==================== 第三章：模型假设 ====================
% BZD 数模社制作——国赛 B 题通用模板
% 使用时请根据题目条件和后续模型逐一替换【】中的内容；假设数量可增减。
%
% ✓ 自查清单：
%   □ 是否只保留后续建模真正需要、且题目没有直接给出的假设？
%   □ 每条假设是否说明依据、简化对象、适用范围及其对模型的作用？
%   □ 是否保留核心机理、物理约束和业务约束，只忽略次要因素？
%   □ 数据假设是否与附件质量检查和实际预处理方法一致？
%   □ 是否避免“数据完全真实可靠”“误差均可忽略”等无依据断言？
%   □ 每条假设是否在后续公式、参数、约束或求解过程中实际使用？
%   □ 若假设不成立，是否能够说明其可能影响和替代处理方式？
%   □ 是否删除重复假设和仅为计算方便而设置的不合理假设？
%%%%%%%%%%%%%%%%%%%%%%%%%%%%%%%%%%%%%%%%%%%%%%%%%%%%%%%%%%%%%

\section{模型假设}

在保留问题核心机理和主要约束的前提下，为明确研究边界并使模型具有可求解性，本文根据题目条件和各问建模需要作出如下必要假设：

\begin{enumerate}
  \setlength{\itemsep}{0.8em}

  \item 假设【具体内容】。

  \textbf{依据与作用：}【说明该假设的题设依据、数据依据、客观规律或合理来源；指出被简化的对象、适用的问题及其对变量、参数、约束或求解过程的作用。】

  \item 假设【具体内容】。

  \textbf{依据与作用：}【说明该假设的题设依据、数据依据、客观规律或合理来源；指出被简化的对象、适用的问题及其对变量、参数、约束或求解过程的作用。】

  \item 假设【具体内容】。

  \textbf{依据与作用：}【说明该假设的题设依据、数据依据、客观规律或合理来源；指出被简化的对象、适用的问题及其对变量、参数、约束或求解过程的作用。】

  \item 假设【具体内容】。

  \textbf{依据与作用：}【说明该假设的题设依据、数据依据、客观规律或合理来源；指出被简化的对象、适用的问题及其对变量、参数、约束或求解过程的作用。】

  \item 假设【具体内容】。

  \textbf{依据与作用：}【说明该假设的题设依据、数据依据、客观规律或合理来源；指出被简化的对象、适用的问题及其对变量、参数、约束或求解过程的作用。】
\end{enumerate}

%%%%%%%%%%%%%%%%%%%%%%%%%%%%%%%%%%%%%%%%%%%%%%%%%%%%%%%%%%%%%
% ==================== 第四章：符号说明 ====================
% BZD 数模社制作——国赛 B 题通用模板
% 使用时请根据全文实际采用的符号替换【】中的内容，并删除多余行。
%
% ✓ 自查清单：
%   □ 是否采用三线表，并按集合/索引、变量、参数、指标的逻辑排列？
%   □ 是否只收录全文反复使用的主要符号？
%   □ 是否删除仅出现一次、可在公式后直接解释的局部符号？
%   □ 有量纲变量是否标明单位，无量纲量是否明确写“无量纲”？
%   □ 同一符号是否只表示一个含义？
%   □ 同一变量是否始终使用同一符号、上下标和单位？
%   □ 是否检查了题面、模型、代码和结果表之间的符号与单位一致性？
%%%%%%%%%%%%%%%%%%%%%%%%%%%%%%%%%%%%%%%%%%%%%%%%%%%%%%%%%%%%%

\section{符号说明}

为便于后续模型表达，现将全文反复使用的主要符号汇总如表~\ref{tab:symbols} 所示。仅在局部公式中出现的符号，将在正文首次出现时另行说明。

\begin{table}[htbp]
  \centering
  \caption{主要符号说明}
  \label{tab:symbols}
  \renewcommand{\arraystretch}{1.25}
  \begin{tabular}{@{}
    >{\centering\arraybackslash}p{0.18\textwidth}
    >{\raggedright\arraybackslash}p{0.56\textwidth}
    >{\centering\arraybackslash}p{0.14\textwidth}@{}}
    \toprule
    \textbf{符号} & \centering\textbf{含义}\arraybackslash & \textbf{单位} \\
    \midrule
    【集合或索引符号】 & 【集合、对象编号、阶段编号或时间索引的含义】 & 【无量纲】 \\
    【状态或观测变量】 & 【描述研究对象状态、位置、数量或观测结果的变量含义】 & 【实际单位】 \\
    【决策变量】 & 【模型需要确定的方案、数量、路径、时刻或控制变量含义】 & 【实际单位】 \\
    【模型参数】 & 【由题目、附件、文献或前序问题确定的参数含义】 & 【实际单位】 \\
    【目标函数或评价指标】 & 【优化目标、误差指标、效益指标或约束度量的含义】 & 【单位或无量纲】 \\
    【其他高频符号】 & 【全文反复使用且无法归入上述类别的符号含义】 & 【单位】 \\
    \bottomrule
  \end{tabular}
\end{table}

\noindent\textbf{注：}填写时应保留符号的上下标、向量或矩阵格式，并确保其与模型公式、程序变量和结果表中的定义一致。

%%%%%%%%%%%%%%%%%%%%%%%%%%%%%%%%%%%%%%%%%%%%%%%%%%%%%%%%%%%%%
% ==================== 第五章：模型的建立与求解 ====================
% BZD数模社制作，本文档为论文模板；如需详细说明，可见论文模板完整版。
% ✓ 使用说明：
%   □ 四个问题均按“思路—建立—求解—结果—检验”组织。
%   □ 数据处理不是固定小节；只有模型确实需要时，才在建模思路或求解步骤中简述。
%   □ 模型建立末尾必须汇总最终模型，不能只分散给出若干公式。
%   □ 模型求解使用 Step 1、Step 2……说明真实求解流程，步骤数量按实际算法调整。
%   □ 结果与分析只引用真实计算结果，禁止预置模型、指标或数值。
%   □ 模型检验与灵敏度分析只选择与本问匹配的 1～2 项，不机械堆砌。
%   □ 预测类可检验误差或残差；评价类可检验排序或权重稳定性；
%     优化类可检验约束可行性、基准对比或重复求解稳定性；
%     仿真与机理类可检验边界条件、守恒关系或参数扰动。
%   □ 数组越界、维度一致、程序成功运行等属于代码自检，不写成模型检验。

\section{模型的建立与求解}

%%%%%%%%%%%%%%%%%%%%%%%%%%%%%%%%%%%%%%%%%%%%%%%%%%%%%%%%%%%%%
% ==================== 5.1：问题一的模型建立与求解 ====================
% ✓ 自查：模型选择有依据；模型已汇总；步骤可复现；结果真实；检验与模型匹配。

\subsection{问题一的模型建立与求解}

\subsubsection{建模思路与模型选择}

问题一要求【概括问题一的任务与目标】，其核心是回答【需要解决的关键数学问题】。
根据【题目条件、研究对象或数据特征】，该问题属于【数学问题类型】。
综合比较【候选方法】在【适用条件、解释能力或求解效率】方面的特点，本文选择
【主模型名称】刻画【变量关系或现实机理】，并以【必要的对比方法】作为参照。
若本问确有数据处理需求，则采用【处理方法】处理【缺失、异常、量纲或口径问题】，理由是【选择依据】；
所得【参数、特征或中间结果】将作为问题二的输入。

\subsubsection{模型的建立}

定义【决策变量、状态变量或解释变量】为【符号及含义】，定义【关键参数】为【符号及含义】。
依据【理论、机理或题目规则】建立【核心方程、目标函数或判别规则】，并结合
【资源限制、边界条件或逻辑关系】给出相应约束：

\begin{equation}
  \text{【核心方程、目标函数或判别规则】},
  \label{eq:q1_core}
\end{equation}

\begin{equation}
  \text{【约束条件、边界条件或适用条件】}.
  \label{eq:q1_constraint}
\end{equation}

\noindent\textbf{模型汇总：}
将上述变量关系、约束条件和参数范围统一整理，得到问题一的完整模型

\begin{equation}
  \mathcal{M}_1:\quad
  \left\{
  \begin{aligned}
    &\text{【核心模型表达式】},\\
    &\text{【约束或适用条件】},\\
    &\text{【变量范围、初始条件或边界条件】}.
  \end{aligned}
  \right.
  \label{eq:q1_summary}
\end{equation}

\subsubsection{模型的求解}

\noindent\textbf{Step 1：}整理【原始输入、题设参数或初始条件】，得到【可直接用于求解的输入形式】。

\noindent\textbf{Step 2：}按照【算法或计算方法】计算【关键变量、中间量或候选解】，并依据【规则】完成【核心求解过程】。

\noindent\textbf{Step 3：}根据【终止条件、筛选准则或评价指标】确定最终结果，并保存【问题二需要调用的输出】。

求解使用【软件、求解器或程序语言】完成，关键参数取值及其来源为【参数设置与依据】，完整程序见附录【编号】。

\subsubsection{求解结果与分析}

求解得到【真实的核心结果】。由【图或表编号】可知，【描述主要现象或变化规律】；
其中【关键结果】说明【数学含义或现实含义】。结合题目要求，问题一的结论为【直接回答问题一】，
并将【真实的中间结果】传递至问题二。

\subsubsection{模型检验与灵敏度分析}

针对【本问模型类型与主要结论】，选取【检验方法】进行验证。
检验指标【指标名称】的结果为【真实结果】，与【判定标准或基准方法】对比后表明【检验结论】。
若【关键参数】具有现实不确定性且会影响主要结论，则在【有依据的变化范围】内进行扰动，
分析【输出数值或决策方案】的变化并说明【稳定区间、临界点或敏感方向】；若无必要，可删除本句而不强行开展灵敏度分析。

%%%%%%%%%%%%%%%%%%%%%%%%%%%%%%%%%%%%%%%%%%%%%%%%%%%%%%%%%%%%%
% ==================== 5.2：问题二的模型建立与求解 ====================
% ✓ 自查：问题一输出已真实进入本问；新增条件已体现；比较标准明确；检验不过度展开。

\subsection{问题二的模型建立与求解}

\subsubsection{建模思路与模型选择}

问题二在问题一的基础上进一步要求【概括问题二的任务】，新增了【新增条件或研究要求】。
本问以问题一得到的【参数、特征或中间结果】作为【本问中的作用】，重点解决【核心数学问题】。
根据【新增条件和数据特点】，选择【主模型名称】对问题一的方法进行【扩展、修正或重新建模】，
并采用【必要的对比模型或基准方案】从【比较指标】方面评价改进效果。
如需额外处理数据，则说明【处理对象、方法及理由】，不单独设置无实际内容的数据预处理小节。

\subsubsection{模型的建立}

在问题一符号体系的基础上，新增【变量或参数】表示【含义】，其余符号与前文保持一致。
依据【理论、机理或题目规则】建立本问的【核心方程、目标函数或判别规则】，并加入【新增约束或条件】：

\begin{equation}
  \text{【问题二核心方程、目标函数或判别规则】},
  \label{eq:q2_core}
\end{equation}

\begin{equation}
  \text{【问题二新增约束、边界条件或适用条件】}.
  \label{eq:q2_constraint}
\end{equation}

\noindent\textbf{模型汇总：}
综合问题一的输入、本问新增条件和变量范围，得到问题二的完整模型

\begin{equation}
  \mathcal{M}_2:\quad
  \left\{
  \begin{aligned}
    &\text{【核心模型表达式】},\\
    &\text{【继承条件与新增约束】},\\
    &\text{【变量范围、初始条件或边界条件】}.
  \end{aligned}
  \right.
  \label{eq:q2_summary}
\end{equation}

\subsubsection{模型的求解}

\noindent\textbf{Step 1：}读取问题一输出的【中间结果】并结合【新增数据或条件】，构造问题二的求解输入。

\noindent\textbf{Step 2：}采用【算法或计算方法】求解【主模型】，记录【目标值、状态量或评价指标】。

\noindent\textbf{Step 3：}在相同输入和评价口径下求解【对比模型或基准方案】，根据【比较指标】确定最终结果并输出【传递给问题三的内容】。

求解使用【软件、求解器或程序语言】完成，关键参数和终止条件为【真实设置及依据】，完整程序见附录【编号】。

\subsubsection{求解结果与分析}

问题二得到【真实的核心结果】。由【图或表编号】可知，【主模型】在【比较指标】方面表现为【真实对比结果】，
说明新增的【条件、机制或改进措施】对【研究对象】产生了【实际影响】。
据此，问题二的答案为【直接结论】，并将【参数、方案或预测结果】用于问题三。

\subsubsection{模型检验与灵敏度分析}

围绕【本问主要结论】，采用【与模型类型匹配的检验方法】验证【拟合、排序、约束或机制】是否可靠。
指标【指标名称】及其真实结果为【结果】，依据【判定标准或基准】可得【检验结论】。
仅当【关键参数】存在现实波动且可能改变结论时，分析其在【有依据的范围】内变化对
【核心输出或最优策略】的影响，并报告【稳定区间、敏感方向或策略切换点】。

%%%%%%%%%%%%%%%%%%%%%%%%%%%%%%%%%%%%%%%%%%%%%%%%%%%%%%%%%%%%%
% ==================== 5.3：问题三的模型建立与求解 ====================
% ✓ 自查：前两问结果与本问条件衔接清楚；模型完整；步骤可复现；结论由真实结果支持。

\subsection{问题三的模型建立与求解}

\subsubsection{建模思路与模型选择}

问题三要求在前两问结果的基础上完成【概括问题三的任务】，需要同时考虑【关键条件一】与【关键条件二】。
本问将【问题一结果】和【问题二结果】分别作为【参数、输入或约束】，其核心是【关键数学问题】。
根据【问题结构、变量性质或求解规模】，选择【主模型名称】描述【目标与约束之间的关系】，
必要时以【基准方法或对比模型】检验所选方案的有效性。
若新增数据必须处理，则在此说明【处理方法及理由】，不固定设置数据预处理标题。

\subsubsection{模型的建立}

定义【本问新增变量】为【符号及含义】，并沿用前两问中的【已有变量或参数】。
依据【现实机理、数学原理或题目规则】建立【核心方程、目标函数或状态转移关系】，同时满足【主要约束】：

\begin{equation}
  \text{【问题三核心方程、目标函数或状态转移关系】},
  \label{eq:q3_core}
\end{equation}

\begin{equation}
  \text{【问题三约束、边界条件或适用条件】}.
  \label{eq:q3_constraint}
\end{equation}

\noindent\textbf{模型汇总：}
汇总前两问输入、本问目标和全部约束，得到问题三的完整模型

\begin{equation}
  \mathcal{M}_3:\quad
  \left\{
  \begin{aligned}
    &\text{【核心模型表达式】},\\
    &\text{【全部主要约束或演化规则】},\\
    &\text{【变量范围、初始条件或边界条件】}.
  \end{aligned}
  \right.
  \label{eq:q3_summary}
\end{equation}

\subsubsection{模型的求解}

\noindent\textbf{Step 1：}整合前两问输出与【本问新增条件】，形成统一的【参数集合、状态初值或候选方案集】。

\noindent\textbf{Step 2：}采用【算法或计算方法】依次完成【关键计算过程】，并记录【必要的中间结果】。

\noindent\textbf{Step 3：}依据【目标、终止条件或筛选准则】获得最终【预测、方案或机理参数】，并输出问题四所需的【结果】。

求解使用【软件、求解器或程序语言】完成，关键参数、初值和终止条件为【真实设置及依据】，完整程序见附录【编号】。

\subsubsection{求解结果与分析}

求解结果为【真实结果】。由【图或表编号】可见，【描述核心规律、方案特征或变化趋势】；
进一步比较【对象或情景】可知【差异及原因】。因此，问题三的结论为【直接回答问题三】，
其中【必要输出】将作为问题四综合分析的依据。

\subsubsection{模型检验与灵敏度分析}

根据【本问模型类型】，采用【检验方法】检查【关键数学关系、约束可行性或结果精度】。
真实检验结果【指标或现象】达到【判定标准或基准】，说明【检验结论】。
若【核心不确定参数】可能影响最终方案，则在【现实变化范围】内进行扰动，重点分析
【数值变化是否引起决策结构变化】，并报告【稳定区间、临界阈值或敏感方向】。

%%%%%%%%%%%%%%%%%%%%%%%%%%%%%%%%%%%%%%%%%%%%%%%%%%%%%%%%%%%%%
% ==================== 5.4：问题四的模型建立与求解 ====================
% ✓ 自查：综合使用前三问成果；最终结论逐项回应题目；检验服务于结论而非形式完整。

\subsection{问题四的模型建立与求解}

\subsubsection{建模思路与模型选择}

问题四要求综合前三问结果完成【最终预测、决策、优化或机理解释任务】。
本问以【问题一结果】、【问题二结果】和【问题三结果】作为【输入、参数或备选方案】，
同时考虑【最终新增条件或现实限制】，需要回答【核心问题】。
据此选择【主模型名称或综合分析方法】完成【最终任务】，并以【评价标准或基准方案】衡量结果质量。
如有必要的数据整理，只说明与最终模型直接相关的【处理内容与理由】。

\subsubsection{模型的建立}

定义【最终决策变量、输出变量或状态变量】为【符号及含义】，定义【新增参数】为【符号及含义】。
结合前三问结论和【现实约束】，建立【核心方程、目标函数、决策规则或机理关系】：

\begin{equation}
  \text{【问题四核心方程、目标函数、决策规则或机理关系】},
  \label{eq:q4_core}
\end{equation}

\begin{equation}
  \text{【问题四约束、边界条件或适用条件】}.
  \label{eq:q4_constraint}
\end{equation}

\noindent\textbf{模型汇总：}
综合全部输入、目标、约束和变量范围，得到问题四的完整模型

\begin{equation}
  \mathcal{M}_4:\quad
  \left\{
  \begin{aligned}
    &\text{【最终模型表达式】},\\
    &\text{【全部主要约束或决策规则】},\\
    &\text{【变量范围、初始条件或边界条件】}.
  \end{aligned}
  \right.
  \label{eq:q4_summary}
\end{equation}

\subsubsection{模型的求解}

\noindent\textbf{Step 1：}汇总前三问的【真实输出】和问题四的【新增输入】，统一符号、单位与评价口径。

\noindent\textbf{Step 2：}采用【算法、求解器或综合分析流程】计算【最终预测、决策变量或候选方案】。

\noindent\textbf{Step 3：}根据【评价标准、可行性条件或决策准则】筛选最终结果，并形成【题目要求的输出形式】。

求解使用【软件、求解器或程序语言】完成，关键参数与终止条件为【真实设置及依据】，完整程序见附录【编号】。

\subsubsection{求解结果与分析}

最终得到【真实的预测值、最优方案、决策建议或机理结论】。
由【图或表编号】可知，【主要规律或方案特征】，其产生原因是【结合模型与现实条件解释原因】。
上述结果分别回答了【题目要求一】与【题目要求二】，最终结论为【完整回答问题四】。

\subsubsection{模型检验与灵敏度分析}

围绕最终结论，选择【与主模型匹配的检验方法】验证【结果精度、约束可行性、稳定性或机理合理性】。
检验得到【真实指标、对照结果或边界表现】，依据【判断标准】可得【检验结论】。
若存在会改变最终决策的【现实不确定参数】，则在【有依据的范围】内分析其影响，
说明【主要趋势、稳定区间或策略切换点】；不得用程序运行正常代替数学意义上的模型检验。


% 可选：仅当全文采用同一模型或同类模型并需要统一检验时启用；
% 启用后应删除各问题中重复的“模型检验与灵敏度分析”小节。
% %%%%%%%%%%%%%%%%%%%%%%%%%%%%%%%%%%%%%%%%%%%%%%%%%%%%%%%%%%%%%
% ==================== 可选：模型的分析与检验 ====================
% 本文件默认不在 main.tex 中启用。
% 仅当全文采用同一模型或同类模型，且需要统一比较检验指标、关键参数或误差传播时启用。
% 启用后必须删除各问题中重复的“模型检验与灵敏度分析”小节，避免重复论证。
% ✓ 自查清单：
%   □ 是否说明被统一检验的模型、指标和判定标准？
%   □ 参数选择及变化范围是否有现实或统计依据？
%   □ 是否分析输入扰动、假设变化、误差来源与传播？
%   □ 所有指标和结论是否来自真实计算记录？

\section{模型的分析与检验}

本文对【需要统一分析的模型或同类模型】进行集中检验。
根据各模型承担的论文任务，分别选取【量化指标】评价【拟合、预测、排序、约束可行性或求解稳定性】，并以
【理论条件、统计阈值、基准方法或题目要求】作为判定标准。真实检验结果见【图表或结果位置】，由此可得【总体检验结论】。

\subsection{灵敏度分析}

选取具有现实不确定性且可能改变主要结论的参数【参数名称】进行灵敏度分析，其基准值和变化范围分别为
【基准值】和【有依据的变化范围】，选择依据是【数据误差、现实波动或参数来源】。
在保持【控制变量】不变的条件下重新求解模型，记录【目标值、关键状态或决策方案】的变化。
结果表明【真实变化规律】，模型在【稳定区间】内保持【结论或方案】不变；当参数达到【临界值】时，
【数值或决策结构】发生【真实变化】，说明【敏感性结论及现实含义】。

\subsection{误差与稳健性分析}

模型误差主要来自【测量、抽样、参数估计、模型简化或数值计算等真实来源】，其可能通过
【变量、方程或求解步骤】传递至【核心输出】。采用【误差指标、输入扰动、假设替换、重复实验或压力测试】
进行分析，得到【真实误差或稳定性结果】。与【允许误差、置信区间或基准方案】比较后可知，
模型在【适用范围】内具有【稳健性结论】；在【边界条件】下仍存在【需要说明的风险或局限】。

%%%%%%%%%%%%%%%%%%%%%%%%%%%%%%%%%%%%%%%%%%%%%%%%%%%%%%%%%%%%%
% ==================== 模型的评价、改进与推广 ====================
% BZD数模社制作，本文档为论文模板；如需详细说明，可见论文模板完整版。
% ✓ 自查清单：
%   □ 是否概括四问模型、输入输出关系和主要结论？
%   □ 每项优点是否对应实际做法或量化证据？
%   □ 每项不足是否写清来源、影响和适用边界？
%   □ 改进措施是否逐项对应前述不足？
%   □ 推广是否说明需要重新标定的变量、参数和约束？
%   □ 是否避免重复前文的模型检验与灵敏度分析？

\section{模型的评价、改进与推广}

本文围绕【研究问题】，针对四个子问题分别采用【各问主要模型或方法】完成
【问题一目标】、【问题二目标】、【问题三目标】和【问题四目标】。
其中，问题一得到的【结果】为问题二提供【输入或参数】，问题二与问题三的【中间结果】
进一步进入问题四的【目标函数、约束或综合决策】。结合前文真实求解结果和检验结论，
现对模型的优点、不足、改进方向及推广价值进行总结。

\subsection{模型的优点}

\begin{enumerate}
  \item 模型具有较强针对性。针对【题目特点】，在【基础模型或理论】中引入【变量、机制或约束】，使模型能够描述【实际关系】，并由【真实结果或指标】提供证据。
  \item 模型结构较为完整。通过【目标函数、状态方程、约束体系或求解策略】统一刻画【核心矛盾】，其中【量化指标或验证结果】达到【真实数值或判定标准】。
  \item 各问题衔接清晰。前序问题输出的【参数、特征或方案】真实进入后续模型，全文符号、单位和评价口径保持一致。
  \item 结果具有【可解释性、稳定性或计算效率】。由【对比实验、误差检验或稳健性分析】可知，【填写有证据支持的具体优点】。
\end{enumerate}

\subsection{模型的不足}

\begin{enumerate}
  \item 受【数据来源、样本范围或统计口径】限制，模型尚未充分考虑【缺失因素】，可能使【具体结果】在【适用边界】之外产生偏差。
  \item 为保证模型可解，对【现实过程】作出【具体假设或简化】，因此模型在【极端工况或特殊情景】下的解释能力受到限制。
  \item 【参数估计、算法规模或随机求解过程】仍存在【具体局限】，可能影响【精度、求解时间或决策稳定性】。
\end{enumerate}

\subsection{模型的改进}

针对上述不足，可从以下方面改进：

\begin{enumerate}
  \item 针对数据局限，补充【数据类型】并扩大【时间、空间或样本】范围，统一【统计口径】，重新估计【相关参数】。
  \item 针对模型简化，引入【新增变量、动态机制或现实约束】，并通过【检验方法】评价改进前后的差异。
  \item 针对计算与不确定性问题，采用【改进算法、鲁棒优化或不确定性建模方法】，比较【精度、稳定性与计算成本】后确定改进方案。
\end{enumerate}

\subsection{模型的推广}

本文模型中的【可迁移结构、变量关系或求解思想】可推广至【相近对象或应用场景】。
推广时应重新收集【新场景数据】，标定【关键参数】，替换【场景相关变量】，并增加
【新的资源、边界或业务约束】；随后采用【验证方法】确认模型在新场景中的适用性，
不得直接照搬本文参数、结果或结论。


\newpage
%%%%%%%%%%%%%%%%%%%%%%%%%%%%%%%%%%%%%%%%%%%%%%%%%%%%%%%%%%%%%
% ==================== AI工具使用声明 ====================
% 2026 国赛要求本声明位于参考文献之前，且两种情形只能保留一种。
% 使用本 Skill、Codex 或其他 AI 工具完成竞赛工作时，必须保持 \cumcmaiusedtrue。
% 只有竞赛全过程确实未使用任何 AI 工具时，才可改为 \cumcmaiusedfalse。

\newif\ifcumcmaiused
\cumcmaiusedtrue

\section*{AI工具使用声明}

\ifcumcmaiused
  本参赛队在竞赛过程中使用了AI工具，主要用于【按实际填写语言润色、代码调试、结果核验、图表优化等简要用途】，详细使用情况见支撑材料。
\else
  本参赛队在竞赛过程中未使用任何AI工具。
\fi

\newpage
\referencescn
\newpage
\appendixcn

\end{document}
